\chapter{Une carte de l'industrie}\label{ch:in:a-map-of-the-industry}

La classification sectorielle du gouvernement des États-Unis comporte une entrée pour «~market making for securities~». Elle la range
sous le code 523150, Investment Banking and Securities Intermediation, à côté des départements de prise ferme des plus grandes
banques. Une firme qui ne négocie que son propre argent est un «~securities speculator for own account~» et relève du code 523910,
Miscellaneous Intermediation, avec «~individuals investing in financial contracts on own account~». De l'autre côté de
la même industrie des titres, le code 6211 de la SEC couvre indifféremment une banque d'investissement, un teneur de marché électronique et une société de
gestion avec 14~billions de dollars d'encours, au même titre. Aucun de ces codes n'indique à qui appartient le capital à risque, qui paie la firme ou combien de temps
elle détient une position, et ces trois réponses déterminent ce que recouvre un poste. Ce chapitre construit la carte dont se sert le reste du
livre, et montre comment situer un employeur sur cette carte à partir des seules sources publiques.

\section{Cinq questions pour situer un employeur}

Le livre~1 a défini les firmes qui négocient~: le teneur de marché, la société de trading pour compte propre, le hedge fund et la
plateforme multigérants, la société de gestion et le desk de la banque. Pour qui choisit où travailler, cinq
questions les distinguent plus finement que n'importe quel nom.

\begin{enumerate}
\item \textbf{À qui appartient le capital à risque~?} Celui des propriétaires (une firme pour compte propre), celui d'investisseurs extérieurs (un fonds), celui des clients sous
un mandat (une société de gestion), celui des actionnaires dans un bilan réglementé (une banque).
\item \textbf{Qui paie la firme~?} Les contreparties par la fourchette, les investisseurs par des frais sur les actifs et sur les profits,
les clients par des commissions, ou personne d'autre que le marché.
\item \textbf{Combien de temps une position est-elle tenue~?} Quelques secondes pour un teneur de marché, de quelques minutes à quelques jours pour une firme à fréquence moyenne,
des mois pour un adepte du suivi de tendance, des années pour un propriétaire d'actifs.
\item \textbf{Que vend-elle, et sur quels marchés~?} Un service (liquidité, exécution, un produit, une licence) dans une ou
plusieurs classes d'actifs.
\item \textbf{Quelle est sa taille~?} En effectifs, en capital et en revenus~: des dizaines, des centaines, des milliers ou des centaines de milliers
de salariés.
\end{enumerate}

\begin{definition}[Typologie des employeurs]\label{def:in:a-map-of-the-industry:taxonomy}
Une \emph{typologie des employeurs}\index{typologie des employeurs} du secteur des marchés classe un employeur, ou une ligne d'activité d'un groupe,
selon son modèle d'affaires -- les réponses aux trois premières questions -- puis selon le produit et la taille. Ce livre
retient treize modèles d'affaires répartis en quatre groupes~: négociation pour compte propre (teneur de marché, firme à fréquence moyenne pour compte propre, société de
trading crypto), gestion d'investissement (fonds systématique, plateforme multigérants, fonds discrétionnaire, société de gestion, propriétaire
d'actifs), banque, et infrastructure (bourse, courtier, fournisseur de données, régulateur).
\end{definition}

\begin{center}\small
\begin{tabular}{@{}>{\raggedright\arraybackslash}p{2.7cm}>{\raggedright\arraybackslash}p{2.5cm}>{\raggedright\arraybackslash}p{4.0cm}>{\raggedright\arraybackslash}p{3.4cm}@{}}
\toprule
modèle d'affaires & capital engagé par & qui paie & durée de détention\\
\midrule
teneur de marché & propriétaires & la fourchette et les rabais & de quelques secondes à quelques heures\\
firme à fréquence moyenne pour compte propre & propriétaires & profit de trading & de quelques minutes à quelques jours\\
fonds systématique & investisseurs & frais sur les actifs et sur les profits & de quelques minutes à quelques mois\\
plateforme multigérants & investisseurs & coûts pass-through et profits & de quelques jours à quelques mois\\
société de gestion & clients & frais sur les actifs & de quelques mois à quelques années\\
propriétaire d'actifs & bénéficiaires & son propre budget & des années\\
division marchés de la banque & actionnaires & fourchette et activité clientèle & de quelques minutes à quelques mois\\
bourse, courtier, fournisseur & actionnaires & frais de transaction, de données et de licence & néant\\
\bottomrule
\end{tabular}
\end{center}

\begin{example}[Trois employeurs, cinq réponses]\label{ex:in:a-map-of-the-industry:three}
Un commodity trading advisor (CTA) qui gère des programmes de suivi de tendance (trend following) sur contrats à terme engage le capital d'investisseurs, est rémunéré par une commission de gestion
et d'une commission de performance, garde des positions de quelques semaines à quelques mois, vend un produit de managed futures sur chaque marché de contrats à terme
liquide, et peut employer quelques dizaines de personnes. Un teneur de marché sur ETF engage le capital de ses propriétaires, est rémunéré par
la fourchette entre le cours du fonds et celui de son panier, garde des positions de quelques minutes à quelques heures, et vend de la liquidité
sur les fonds qu'il cote. Le desk de flux de taux d'une banque engage le capital des actionnaires dans un bilan réglementé, est rémunéré
par la fourchette et la relation client, garde ses positions pendant quelques jours et fournit des cours à des clients institutionnels. Trois
postes appelés «~trader~» qui n'ont en commun que les écrans.
\end{example}

\begin{remark}[Groupes et entités juridiques]\label{rem:in:a-map-of-the-industry:groups}
Un groupe peut posséder un teneur de marché, un fonds et une filiale technologique sous une même marque, et la division marchés d'une banque
regroupe des activités de plusieurs natures. La typologie classe l'entité juridique ou la ligne d'activité que rejoint un salarié,
et non la marque. Lorsque les activités d'un groupe diffèrent, ses offres d'emploi diffèrent aussi.
\end{remark}

\section{La typologie~: modèles d'affaires et leurs sous-types}

La \Cref{fig:in:a-map-of-the-industry:tree} représente les treize modèles dans leurs quatre groupes, avec les chapitres de ce livre
qui décrivent chacun en tant qu'employeur. Le groupe de la négociation pour compte propre est celui où les associations professionnelles des bourses tracent
leur propre frontière.

\begin{definition}[Firme de trading pour compte propre]\label{def:in:a-map-of-the-industry:ptf}
Une \emph{firme de trading pour compte propre}\index{firme de trading pour compte propre} (principal trading firm) est une société de trading pour compte propre (livre~1) au sens
que retiennent les associations professionnelles des bourses~: une firme qui ne négocie que son propre capital, sur les bourses et les plateformes
de négociation électroniques, surtout en tant que fournisseur de liquidité, et n'a aucun client dont elle investit l'argent.
\end{definition}

L'association américaine réserve l'adhésion de son groupe de négociants pour compte propre aux «~firms trading their own
capital~»~; l'association européenne décrit ses membres comme des «~independent market makers and providers of liquidity
and risk transfer~». Ce terme compte pour un salarié, car il fixe deux choses à la fois~: chaque dollar de revenu est
un revenu de trading, et personne d'autre que les propriétaires n'a de droit sur le profit (livre~16, chapitre~2).

\begin{omfigure}
\begin{tikzpicture}[font=\footnotesize,>=stealth,
  grp/.style={draw,rounded corners,fill=omDef!12,minimum width=2.9cm,minimum height=0.6cm,align=center},
  mdl/.style={draw,rounded corners,fill=omProp!8,minimum width=5.9cm,minimum height=0.42cm,align=center,inner sep=2pt}]
\foreach \m/\y/\c in {a/0/{teneur de marché (ch.~2, 3)},b/-0.55/{firme à moyenne fréquence (ch.~6)},c/-1.1/{firme de trading crypto (ch.~10)},
  d/-1.9/{fonds systématique (ch.~4)},e/-2.45/{plateforme multigérants (ch.~5)},f/-3.0/{fonds discrétionnaire (ch.~4, 5)},
  g/-3.55/{gestionnaire d'actifs (ch.~8)},h/-4.1/{propriétaire d'actifs (ch.~8)},i/-4.9/{division marchés d'une banque (ch.~7)},
  j/-5.7/{bourse (ch.~9)},k/-6.25/{courtier (ch.~9)},l/-6.8/{fournisseur (ch.~9)},n/-7.35/{régulateur (ch.~9)}}
  {\node[mdl] (\m) at (6.2,\y) {\c};}
\node[grp] (P) at (0,-0.55) {négociation pour compte propre};
\node[grp] (I) at (0,-3.0) {gestion d'investissement};
\node[grp] (B) at (0,-4.9) {banque};
\node[grp] (N) at (0,-6.525) {infrastructure};
\foreach \g/\ms in {P/{a,b,c},I/{d,e,f,g,h},B/{i},N/{j,k,l,n}} {\foreach \m in \ms {\draw[->,gray] (\g.east) -- (\m.west);}}
\node[font=\scriptsize,align=center] at (0,0.55) {groupe};
\node[font=\scriptsize,align=center] at (6.2,0.55) {modèle d'affaires (là où ce livre le décrit)};
\end{tikzpicture}

\omcaption{La \omterm{def:in:a-map-of-the-industry:taxonomy}{typologie des employeurs} de ce livre~: treize modèles d'affaires répartis en quatre groupes, et les chapitres de la partie~I qui
décrivent chacun en tant qu'employeur. Les rôles qui traversent ces modèles relèvent de la partie~III.}\label{fig:in:a-map-of-the-industry:tree}
\end{omfigure}

Au sein de chaque modèle, les sous-types diffèrent assez pour changer le métier. Parmi les teneurs de marché, une firme delta-one électronique qui cote
actions, contrats à terme et devises (chapitre~2) n'est pas une maison d'options qui valorise des milliers de strikes (chapitre~3). Parmi
les sociétés de gestion, un fonds multistratégie piloté par la recherche, un suiveur de tendance et une maison d'arbitrage statistique partagent une
méthode mais pas un rythme (chapitre~4). Les banques se distinguent par la force de leurs produits~: certaines tirent l'essentiel de leurs revenus de marché en dérivés sur actions,
la plupart en taux et en crédit (chapitre~7). La carte est aussi une carte de cette série. Le tableau ci-dessous indique, pour chaque modèle d'affaires, les livres qu'un nouveau venu doit lire.

\begin{center}\small
\begin{tabular}{@{}lp{7.2cm}@{}}
\toprule
modèle d'affaires & Livres de la série qui enseignent ce métier\\
\midrule
teneur de marché & 1 (marchés), 10 (microstructure), 11 (market making), 13--14 (systèmes, réseaux)\\
\omterm{def:in:a-map-of-the-industry:ptf}{firme de trading pour compte propre} à moyenne fréquence & 7 (métier de la recherche), 8 (stratégies), 12 (apprentissage automatique)\\
société de trading crypto & 3 (marchés crypto), 11, 14 (connectivité crypto)\\
fonds systématique & 4 (méthodes), 7, 8, 9 (stratégies), 12\\
plateforme multigérants & 8, 16 (le desk et la firme)\\
société de gestion et propriétaire d'actifs & 1, 8 (les stratégies des sociétés de gestion)\\
division marchés de la banque & 2 (taux, change, crédit), 5 (dérivés), 6 (XVA et risque), 9 (desks bancaires)\\
bourse, courtier, fournisseur, régulateur & 1, 10, 14, 15 (plateformes), 16\\
\bottomrule
\end{tabular}
\end{center}

\section{Taille~: effectifs, capital et revenus}

La taille est la cinquième question et la plus visible. Elle couvre cinq ordres de grandeur, des sociétés de trading comptant
une douzaine de personnes aux banques qui en comptent des centaines de milliers, et elle change le métier plus qu'il n'y paraît~: une firme de cinquante personnes ne dispose
pas d'un département des risques distinct, une firme de cinquante mille en compte plusieurs.

\begin{dated}{2026-09}{Taille des employeurs}\label{dat:in:a-map-of-the-industry:sizes}
\textbf{Effectifs issus des dépôts et des pages des firmes} (fin 2025 sauf indication contraire)~: JPMorganChase 318\,512 employés~;
Goldman Sachs un effectif de 47\,400~; BlackRock environ 24\,900~; CME Group environ 3\,875~; Virtu Financial environ 1\,027
(février 2026). Deux plateformes multigérants indiquent sur leur page d'accueil «~7\,000+ collaborateurs dans le monde~» et «~360+
équipes d'investissement~» (Millennium) et «~2000+ professionnels des équipes d'investissement et du support~» (Balyasny).
\textbf{Décompte des firmes} (FINRA, fin 2024)~: 3\,249 broker-dealers enregistrés, dont 2\,891 comptent 1--150 représentants
enregistrés, 209 en comptent 151--499 et 149 en comptent 500 ou plus~; 32\,090 firmes ne sont enregistrées qu'en tant que conseillers en investissement.
La FINRA classe 103 broker-dealers dans ses segments d'activité de trading pour compte propre et de market making, dont 93 sont petits selon sa
définition, et aucun n'est grand.
\end{dated}

Deux enseignements en découlent. Le nombre de conseillers en investissement enregistrés est dix fois celui des broker-dealers, et il
a progressé alors que celui des broker-dealers a reculé (\cref{fig:in:a-map-of-the-industry:reg})~: la plupart des firmes du secteur
gèrent de l'argent, et la plupart sont petites. Et les \omterm{def:in:a-map-of-the-industry:ptf}{firmes de trading pour compte propre} sont peu nombreuses et surtout petites selon le décompte du régulateur,
qui mesure les représentants enregistrés, non les effectifs~: un teneur de marché qui emploie un millier d'ingénieurs et de traders peut
n'enregistrer que les quelques personnes qui traitent avec le public, si bien que sa taille selon la FINRA ne dit rien de sa taille réelle. Toute taille a une unité, et
l'unité doit être lue.

\begin{omfigure}
\begin{tikzpicture}
\begin{axis}[width=11cm,height=5.4cm,xlabel={\small fin d'exercice},ylabel={\small firmes (milliers)},xmin=2014.6,xmax=2024.4,
  ymin=0,ymax=36,xtick={2015,2016,2017,2018,2019,2020,2021,2022,2023,2024},xticklabel style={/pgf/number format/1000 sep={}},
  tick label style={font=\footnotesize},grid=major,grid style={gray!20},table/col sep=comma,
  legend style={legend cell align=left,font=\scriptsize,at={(0.5,-0.3)},anchor=north,/tikz/every even column/.append style={column sep=8pt}},legend columns=2]
\addplot[omDef,thick,mark=*] table[x=year,y=ia_only]{figdata/industry/01-a-map-of-the-industry/registrations.csv};
\addplot[omThm,thick,mark=square*] table[x=year,y=bd]{figdata/industry/01-a-map-of-the-industry/registrations.csv};
\legend{enregistrées uniquement comme conseillers en investissement,broker-dealers enregistrés auprès de la FINRA}
\end{axis}
\end{tikzpicture}

\omcaption{Firmes enregistrées dans l'industrie américaine des valeurs mobilières, 2015--2024~: broker-dealers (y compris ceux également enregistrés
comme conseillers) sont passés de 3.94 à 3.25 milliers, et les firmes enregistrées uniquement comme conseillers en investissement sont passées de 28.7 à 32.1
milliers. Données~: FINRA Industry Snapshot 2025, tableau 2.1.5, via \texttt{in\_map.finra}.}\label{fig:in:a-map-of-the-industry:reg}
\end{omfigure}

\begin{remark}[Tranches de taille dans ce livre]\label{rem:in:a-map-of-the-industry:bands}
Lorsque ce livre classe les employeurs par taille, il retient les effectifs et non les représentants enregistrés, en quatre tranches~: moins de 50, de 50 à
500, de 500 à 5\,000, et plus de 5\,000. Un effectif est une fourchette datée, jamais un chiffre ponctuel~: les firmes arrondissent, comptent les prestataires autrement,
et changent de taille plus vite que ne le font leurs dépôts.
\end{remark}

\section{Lire un employeur de l'extérieur}

Un observateur extérieur peut situer un employeur à partir de trois types de données publiques~: les codes sectoriels que les agences statistiques et
les régulateurs attribuent, les enregistrements et les rapports périodiques qu'exige la loi, et les publications propres de la firme.

\begin{definition}[Code de classification sectorielle]\label{def:in:a-map-of-the-industry:code}
Un \emph{code de classification sectorielle}\index{code de classification sectorielle} affecte un établissement ou une société à une
branche d'une classification statistique selon son activité principale~: le North American Industry Classification System
(NAICS, à six chiffres), la Standard Industrial Classification qu'utilise encore la SEC sur EDGAR (SIC, à quatre chiffres), ou la
NACE de l'Union européenne.
\end{definition}

Les codes sont grossiers là où le secteur est fin. Le NAICS 2022 a fusionné «~banque d'investissement et négociation de titres~» et
«~courtage de titres~» en un seul code, 523150, et «~gestion de portefeuille~» et «~conseil en investissement~» en 523940~: les
codes qui permettaient de distinguer un dealer d'un courtier ne le permettent plus. La SEC n'attribue un code SIC qu'aux émetteurs de
titres enregistrés, si bien qu'une firme privée sur EDGAR n'en a aucun, quelle que soit sa taille. Les codes sont la première lecture, jamais la dernière.

Les enregistrements en disent davantage. Un broker-dealer dépose un rapport annuel audité (formulaire X-17A-5) et des rapports financiers périodiques
(FOCUS), dont seule la situation financière auditée est publique~; un gestionnaire institutionnel détenant plus de 100~millions de dollars d'actions américaines dépose chaque trimestre ses positions sur le
formulaire 13F~; un émetteur dépose des rapports annuels et trimestriels~; un opérateur de système de négociation alternatif dépose le formulaire ATS-N.
Ce sont des obligations, non des choix, et elles suivent l'activité et non la marque~: chacune renseigne sur ce que l'entité
fait. L'identifiant d'entité juridique (LEI, livre~2, chapitre~28) relie entre elles les entités d'un groupe, et ses données de référence sont
librement réutilisables.

\begin{method}[Situer un employeur à partir des registres publics]\label{met:in:a-map-of-the-industry:place}
\begin{enumerate}
\item Lister les entités juridiques du groupe~: les enregistrements LEI, la recherche par entité d'EDGAR, les registres nationaux des sociétés.
\item Lire les codes sectoriels en connaissant leurs limites~: SIC seulement pour les émetteurs, NAICS grossier, NACE plus grossier encore.
\item Lister les enregistrements de chaque entité et les familles de formulaires qu'elle dépose~: rapports des broker-dealers, rapports de positions,
rapports d'émetteurs, formulaires des opérateurs de plateformes de négociation.
\item Lire la présentation que la firme fait elle-même de son activité, sur son site et dans tout dépôt qui décrit l'activité.
\item Déterminer le modèle d'affaires de l'entité à laquelle appartient le poste, et noter quel registre l'a établi.
\end{enumerate}
\end{method}

\begin{example}[Une firme, trois codes]\label{ex:in:a-map-of-the-industry:codes}
Virtu Financial, la société holding cotée d'un teneur de marché électronique, porte le code SIC 6211, Security Brokers,
Dealers \& Flotation Companies, sur EDGAR. Sa filiale opérationnelle américaine, un broker-dealer, n'a aucun code SIC mais dépose
des rapports annuels de broker-dealer et un formulaire de système de négociation alternatif. Le même code SIC 6211 est porté par Goldman Sachs
et par BlackRock, une société holding bancaire et une société de gestion. Le code range trois employeurs différents dans une même case~; les
formulaires placent la filiale du teneur de marché dans la bonne.
\end{example}

\section{Tutoriel~: classer les employeurs à partir de leurs dépôts}\label{tut:in:a-map-of-the-industry:tut}

\begin{tutorial}
\textbf{Objectif.} Mesurer dans quelle mesure les codes publics et les formulaires déposés par une entité restituent son modèle d'affaires. \textbf{État
final~:} la table des scores ci-dessous et la \cref{fig:in:a-map-of-the-industry:acc}.
\begin{tutsteps}
\item \textbf{L'échantillon.} \texttt{data/industry/employer\_sample.csv} regroupe 21 entités qui déposent sur EDGAR~: onze
entités de market making, quatre conseillers de fonds systématiques, deux plateformes multigérants, deux banques, une société de gestion et une
bourse. Pour chacune, le code SIC et les familles de formulaires présents dans son historique de dépôts récent proviennent des données de soumission
d'EDGAR~; le modèle d'affaires provient de la déclaration propre de la firme, avec une ligne de registre par entité.
\item \textbf{Les règles.} \texttt{firm.industrymap} implémente trois règles
(\cref{lst:in:a-map-of-the-industry:rules})~: le code SIC seul~; les formulaires seuls (un broker-dealer qui n'est pas un émetteur est
une \omterm{def:in:a-map-of-the-industry:ptf}{firme de trading pour compte propre}~; un déposant de positions 13F qui n'est ni l'un ni l'autre est un gestionnaire d'investissement)~; et les formulaires
d'abord, puis le SIC pour les émetteurs.
\omcode{code/firm/industrymap/firm_industrymap.py}{84}{117}{Trois règles qui classent une entité à partir de ses
registres publics.}\label{lst:in:a-map-of-the-industry:rules}
\item \textbf{Score.} \texttt{in\_map.accuracy()} évalue chaque règle à deux niveaux~: le groupe grossier (trading pour compte propre,
gestion d'investissement, banque, infrastructure) et le modèle d'affaires fin.
\item \textbf{Lire les erreurs.} \texttt{in\_map.misses()} liste les entités que la meilleure règle classe mal.
\end{tutsteps}
\end{tutorial}

\begin{center}\small
\begin{tabular}{@{}lrr@{}}
\toprule
règle & bon groupe & bon modèle\\
\midrule
code SIC seul & 3 sur 21 & 3 sur 21\\
formulaires seuls & 15 sur 21 & 13 sur 21\\
formulaires, puis SIC pour les émetteurs & 18 sur 21 & 16 sur 21\\
\bottomrule
\end{tabular}
\end{center}

\begin{omfigure}
\begin{tikzpicture}
\begin{axis}[width=10.5cm,height=5cm,xbar,area legend,bar width=7pt,xmin=0,xmax=100,xlabel={\small part de l'échantillon classée correctement (\%)},
  ytick={0,1,2},yticklabels={code SIC seul,formulaires seuls,{formulaires, puis SIC}},y dir=reverse,ymin=-0.6,ymax=2.6,
  tick label style={font=\footnotesize},grid=major,grid style={gray!20},table/col sep=comma,
  legend style={legend cell align=left,font=\scriptsize,at={(0.5,-0.32)},anchor=north,/tikz/every even column/.append style={column sep=8pt}},legend columns=2]
\addplot[fill=omDef!55,draw=omDef] table[y=k,x=coarse]{figdata/industry/01-a-map-of-the-industry/accuracy.csv};
\addplot[fill=omThm!55,draw=omThm] table[y=k,x=fine]{figdata/industry/01-a-map-of-the-industry/accuracy.csv};
\legend{bon groupe,bon modèle d'affaires}
\end{axis}
\end{tikzpicture}

\omcaption{Classer 21 entités EDGAR à partir de leurs registres publics~: le code SIC restitue 14~\% d'entre elles, les formulaires
qu'elles déposent 71~\% des groupes, et les formulaires avec le SIC pour les émetteurs 86~\%. Données~: soumissions EDGAR et les
déclarations propres des firmes, via \texttt{in\_map.accuracy}.}\label{fig:in:a-map-of-the-industry:acc}
\end{omfigure}

Le code SIC ne place que les cinq émetteurs, et deux d'entre eux à tort~: une banque d'investissement et une société de gestion se retrouvent
sous le code des teneurs de marché. Les formulaires placent chaque teneur de marché privé dont le broker-dealer porte le nom de la firme. Les trois
erreurs de la meilleure règle sont instructives. L'entité d'une firme de trading quantitatif de l'échantillon ne dépose que des déclarations
de positions~: l'entité retenue n'était pas un broker-dealer, et la règle y a vu un gérant de fonds. Et le code SIC de l'émetteur
place mal la banque et la société de gestion. Chaque erreur marque une étape~1 de la méthode, le recensement des entités du groupe, qui a été
sautée.

\textbf{Pour aller plus loin.} Ajouter une règle qui lit un émetteur au code SIC 6211 et ayant des déclarations de positions comme une société de gestion, et
regarder ce qu'elle corrige et ce qu'elle casse (exercice~7)~; ajouter dix entités de son choix, et se demander si ce choix flatte les règles
(exercice~8).

\section{Construction~: la carte des employeurs}\label{bld:in:a-map-of-the-industry:industrymap}

\begin{build}
\bfield{Rôle} Une seule taxonomie des employeurs pour tout le livre, afin que les données de rémunération (chapitre~14), les profils de firmes (chapitres
2--10) et les rôles (partie~III) soient classés de la même façon.
\bfield{Interface} \texttt{firm.industrymap}~: \texttt{MODELS} (modèle d'affaires vers groupe, dont le capital, qui paie, la
période de détention)~; \texttt{GROUP}, \texttt{BOOKS}~; \texttt{Entity(name, cik, sic, families, truth)}, \texttt{load\_sample(path)}~;
\texttt{sic\_class}, \texttt{forms\_class}, et les règles \texttt{by\_sic}, \texttt{by\_forms}, \texttt{combined}~;
\texttt{score(entities, rule, level)} et \texttt{confusion(entities, rule, level)}.
\bfield{Règles} Une règle répond par un modèle d'affaires ou \texttt{unknown}, jamais par une supposition sans registre derrière elle~; chaque
modèle d'affaires a un groupe et une liste de livres~; un score indique son niveau.
\bfield{Tests d'acceptation} \texttt{code/firm/industrymap/tests/}~: chaque modèle a un groupe et des livres~; les règles SIC et formulaires sur
des entités construites~; la règle combinée se replie sur le SIC pour les émetteurs~; les scores aux deux niveaux diffèrent là où la règle
ne sait pas séparer une plateforme d'un fonds systématique.
\bfield{Pour aller plus loin} Ajouter une règle sur le code NAICS qu'un employeur déclare dans ses dépôts de conditions de travail (chapitre~14), et
une règle sur les agréments d'une firme britannique au registre de la FCA.
\end{build}

\begin{omsources}
\item US Census Bureau, North American Industry Classification System 2022~: codes, fichier d'index et table de
correspondance 2017-2022.
\item Données de soumission SEC EDGAR (\texttt{data.sec.gov}) pour les 21 entités de l'échantillon~; Forms 10-K 2025 de
JPMorgan Chase, Goldman Sachs, BlackRock, CME Group et Virtu Financial.
\item Les pages d'accueil des firmes pour leurs descriptions de leur activité (lignes de registre F4--F19).
\item FINRA, Industry Snapshot 2025, tableaux 2.1.3, 2.1.5 et 2.6.1--2.6.3.
\item FIA, pages Principal Traders Group et European Principal Traders Association~; GLEIF, conditions d'utilisation des données LEI.
\end{omsources}

\section{Exercices}

\begin{exercise}[$\star$]\label{exo:in:a-map-of-the-industry:1}
Donner le code NAICS 2022 d'une firme qui fait du market making sur des titres, d'une firme qui ne fait que du trading pour son propre compte, et
d'un conseiller en investissement, puis dites lesquels de ces codes auraient séparé les codes 2017.
\end{exercise}

\begin{exercise}[$\star$]\label{exo:in:a-map-of-the-industry:2}
D'après l'encadré daté, quelle part des broker-dealers enregistrés auprès de la FINRA à la fin de 2024 avait 150 représentants inscrits ou
moins, et quelle part des segments de trading pour compte propre et de market making en avait-elle~?
\end{exercise}

\begin{exercise}[$\star$]\label{exo:in:a-map-of-the-industry:3}
Répondez aux cinq questions de la section~1 pour le fonds indiciel d'un émetteur d'ETF, pour l'équipe actions interne d'un fonds de pension et une
bourse de crypto-actifs.
\end{exercise}

\begin{exercise}[$\star\star$]\label{exo:in:a-map-of-the-industry:4}
La règle SIC classe correctement 3 des 21 entités. Combien en laisse-t-elle en unknown, et pourquoi~?
\end{exercise}

\begin{exercise}[$\star\star$]\label{exo:in:a-map-of-the-industry:5}
De combien, en pourcentage, le nombre de broker-dealers enregistrés auprès de la FINRA a-t-il baissé entre 2015 et 2024, et
de combien le nombre de firmes enregistrées uniquement comme conseillers en investissement a-t-il augmenté, en pourcentage~?
\end{exercise}

\begin{exercise}[$\star\star$]\label{exo:in:a-map-of-the-industry:6}
Dans la \cref{fig:in:a-map-of-the-industry:acc}, pourquoi la barre «~bon modèle~» est-elle plus basse que la barre «~bon groupe~» pour la
règle des formulaires, et quelles deux entités font la différence~?
\end{exercise}

\begin{exercise}[$\star\star\star$]\label{exo:in:a-map-of-the-industry:7}
\emph{Programmation.} Étendez la règle combinée de sorte qu'un émetteur de code SIC 6211 qui dépose aussi des déclarations de positions soit classé comme
société de gestion. Combien des 21 entités place-t-elle désormais dans le bon groupe, et quelle entité la nouvelle règle
classe-t-elle encore à tort, et pourquoi~?
\end{exercise}

\begin{exercise}[$\star\star\star$]\label{exo:in:a-map-of-the-industry:8}
\emph{Trouver la faille.} «~La règle des formulaires place 86~\% des employeurs dans le bon groupe, donc un candidat à l'emploi
peut classer n'importe quelle firme de trading à partir de ses dépôts EDGAR.~»
\end{exercise}

\section{Problème~: classer de l'extérieur}

\begin{problem}[{Problème du week-end --- classer de l'extérieur}]\label{pb:in:a-map-of-the-industry:1}
Un conseiller d'orientation veut un outil qui indique aux étudiants quel type de firme leur a proposé un emploi, à partir de registres publics uniquement.
On dispose de l'échantillon du chapitre et de ses règles.

\medskip\noindent\textbf{Partie I --- Les questions.}
\begin{enumerate}
\item Énoncer les cinq questions qui situent un employeur, et préciser lesquelles, au nombre de trois, définissent son modèle d'activité.
\item Y répondre pour un teneur de marché et pour un fonds systématique.
\item Définir une \omterm{def:in:a-map-of-the-industry:ptf}{firme de trading pour compte propre} (principal trading firm), et préciser ce que les règles d'adhésion des deux associations
ajoutent à la société de trading pour compte propre du livre~1.
\item Pourquoi le chapitre classe-t-il des entités juridiques ou des lignes d'activité plutôt que des marques~?
\item Donner les quatre groupes de la typologie et le nombre de modèles d'activité dans chacun.
\end{enumerate}

\medskip\noindent\textbf{Partie II --- Les codes.}
\begin{enumerate}[resume]
\item Quel code NAICS 2022 couvre le market making de titres, et quel code couvre les spéculateurs pour compte propre~?
\item Quels codes de 2017 ont été fusionnés dans 523150 et dans 523940~?
\item Pourquoi le broker-dealer d'un teneur de marché privé n'a-t-il aucun code SIC sur EDGAR~?
\item Quelles trois entités de l'échantillon portent le code SIC 6211, et quels sont leurs modèles d'activité~?
\item Combien d'entités de l'échantillon la règle SIC place-t-elle correctement, et combien laisse-t-elle inconnues~?
\end{enumerate}

\medskip\noindent\textbf{Partie III --- Les formulaires.}
\begin{enumerate}[resume]
\item Nommer les familles de formulaires que lisent les règles, et ce que chacune établit sur l'entité.
\item Combien d'entités la règle des formulaires place-t-elle dans le bon groupe, et dans le bon modèle d'activité~?
\item Pourquoi la règle des formulaires ne peut-elle pas distinguer une plateforme multigérants d'un fonds systématique~?
\item Combien d'entités la règle combinée place-t-elle correctement, à chaque niveau~?
\item Nommer les trois erreurs de la règle combinée au niveau du groupe et la cause de chacune.
\end{enumerate}

\medskip\noindent\textbf{Partie IV --- Le verdict.}
\begin{enumerate}[resume]
\item Combien de broker-dealers la FINRA a-t-elle comptés à la fin de 2024, et combien dans ses
segments de trading pour compte propre et de market making~?
\item Pourquoi une tranche de taille de la FINRA est-elle une mauvaise mesure de l'effectif d'un teneur de marché~?
\item Énoncer le \emph{résultat nommé}~: la part de l'échantillon que chaque règle classe correctement, au niveau du groupe.
\item Quelle étape de la \cref{met:in:a-map-of-the-industry:place} aurait détecté chacune des erreurs~?
\item En deux phrases, que doit indiquer l'outil du conseiller en plus de sa réponse~?
\end{enumerate}
\end{problem}

\section{Questions d'entretien}

\begin{interviewq}[$\star$ \iqroles{trader, chercheur}]\label{iq:in:a-map-of-the-industry:1}
Qu'est-ce qui distingue un teneur de marché d'un hedge fund~? Donnez trois différences qui comptent pour quelqu'un qui y travaille.
\end{interviewq}

\begin{interviewq}[$\star$ \iqroles{trader}]\label{iq:in:a-map-of-the-industry:2}
Pourquoi cela importe-t-il à un trader qui paie la firme~?
\end{interviewq}

\begin{interviewq}[$\star\star$ \iqroles{développeur, risque}]\label{iq:in:a-map-of-the-industry:3}
Une entité dépose des rapports annuels X-17A-5 auprès de la SEC. Qu'est-ce que cela vous apprend sur elle, et qu'est-ce que cela ne vous apprend pas~?
\end{interviewq}

\begin{interviewq}[$\star\star$ \iqroles{chercheur}]\label{iq:in:a-map-of-the-industry:4}
Estimez le nombre de personnes qui travaillent pour des \omterm{def:in:a-map-of-the-industry:ptf}{firmes de trading pour compte propre} aux États-Unis. Indiquez votre méthode et votre fourchette.
\end{interviewq}

\begin{interviewq}[$\star\star$ \iqroles{banque, risque}]\label{iq:in:a-map-of-the-industry:5}
Le flow desk d'une banque et un teneur de marché détiennent tous deux de l'inventaire pour servir acheteurs et vendeurs. Dans chaque
cas, qui porte le capital exposé au risque, et qu'est-ce qui limite ce que chacun peut détenir~?
\end{interviewq}

\begin{interviewq}[$\star\star\star$ \iqroles{chercheur, ingénieur ML}]\label{iq:in:a-map-of-the-industry:6}
Vous construisez un classifieur qui étiquette les modèles d'activité des employeurs à partir de leurs dépôts publics, et il atteint
86~\% sur un échantillon que vous avez constitué. Comment sauriez-vous s'il faut faire confiance à ce chiffre~?
\end{interviewq}
